\section{Related Work}
\label{sec:related}

\paragraph{Figure-associated text.} \citet{ramesh2015figsum} introduced \figsum{}, a system that summarises a
figure from the text of its article, together with two human gold standards. The first, GS1, lists every
sentence that annotators judged to be associated with each figure, whether or not it names the figure; this is
exactly our task. They compare five unsupervised methods on GS1 (described in Section~\ref{sec:figsumbase});
the best, which simply returns every paragraph that cites the figure, reaches F1 0.66. We found no later work
that reports results on GS1, so this is the state of the art for the task as far as we can establish.
Earlier, \citet{bockhorst2012beyond} linked sentences of the \emph{abstract} to figures; the setting is
narrower (abstract sentences only), but we borrowed their observation that the order of sentences and figures
is informative.

\paragraph{Datasets of figures and their mentions.} Several large datasets pair figures with the text that
cites them. MedICaT \citep{subramanian2020medicat} contains about 217,000 figures from 131,000 open-access
biomedical papers and, for 74\% of the figures, the inline references taken from S2ORC \citep{lo2020s2orc}.
SciCap+ \citep{yang2023scicapplus} extends a figure-captioning dataset with the paragraphs that mention each
figure and with text read from the image by OCR. In both, the associated text is found through explicit
mentions; neither labels implicit references. They are therefore not training data for our task, but they
suggested two ideas that we test: mining free labels from citations (Section~\ref{sec:fakelabels}) and giving
the model text extracted from the figure image (Section~\ref{sec:vis}).

\paragraph{How much does the image help?} Work on generating figure captions has repeatedly found that the
text around a figure carries most of the usable information. \citet{huang2023summaries} show that more than 75\% of
caption words can be aligned with the paragraphs that mention the figure, and that summarising those paragraphs
as text outperforms earlier vision-based captioning. In SciCap+, the figure image alone gives very low scores,
and removing generic image features from the full model does not hurt; the authors describe those features as
acting like noise because the image encoder was not trained on scientific figures. OCR text, in contrast, did
help in their experiments. We return to this comparison when we report our own experiment with figure text.

\paragraph{Document parsing.} GROBID \citep{lopez2009grobid} converts scholarly PDFs to structured XML and is
the standard open tool for this. We use it as one of our parsing routes and first verified our installation
against its published benchmark (figure-reference F1 75.1 in our run, 74.26 published).

\paragraph{Encoders.} Our learned models fine-tune two pre-trained transformer encoders. DeBERTa-v3
\citep{he2021deberta,he2023debertav3} represents the content and the position of each token separately and is
pre-trained by detecting replaced tokens; it is among the strongest encoders for judging pairs of texts.
ModernBERT \citep{warner2024modernbert} is a recent encoder with a context of 8{,}192 tokens, which makes it
possible to read a long stretch of a paper in one pass.
